\section{Incidence-indexed sparse Alexander elimination}
\label{sec:sparse}

The modular Alexander filter is polynomial-time, so improving it does not
by itself repair the recognizer's exponential worst case.  It nevertheless
matters for two reasons.  Large diagrams can be rejected without reaching
homology, and dense storage can consume resources before a favorable
decomposition or invariant has an opportunity to help.  The original
Alexander rows are already sparse; the archived filter expands the chosen
minor to a dense matrix.  This section removes that expansion for large
inputs and analyzes the resulting work.

\subsection{Preserving the existing one-sided witness}

At a crossing, the Wirtinger--Alexander relation contributes at most three
column terms.  Terms in the same column are combined.  Deleting the chosen
row and column gives an $(n-1)\times(n-1)$ minor with $O(n)$ stored polynomial
entries.  The implementation evaluates exactly the same minor as before at
$t=-1$ and at the repository's fixed generic field element.  It also retains
the same allowable unit values $\{\pm t^k:0\leq k<n\}$.

Consequently, the new backend changes neither the mathematical filter nor
its evidence dictionary.  It changes how a numerical determinant is
computed.  A determinant outside the unit set still proves knottedness;
agreement remains inconclusive.  The implementation does not infer a
random-error probability from using a large fixed constant.  In particular,
no independence assumption about the input diagram and the evaluation point
is silently introduced.

The dense routine remains available as a reference and for small matrices.
The default \code{auto} policy selects sparse elimination at $n\geq256$
crossings and dense elimination below that threshold.  Explicit
\code{dense} and \code{sparse} choices support comparison.  This crossover
is an engineering choice informed by pilot measurements, not a universal
theorem about which routine is faster.

\subsection{Pivoting on actual incidences}

Store every active row as a map from column index to a nonzero field value.
For each column also store the set of rows containing it.  A lazy priority
queue chooses a shortest active row.  Among that row's nonzero columns,
choose one having minimum active column degree, with stable index tie
breaking.  This is a local fill heuristic related to the classical
Markowitz approach~\cite{markowitz}; it does not claim to find a global
minimum-fill pivot or even the minimum Markowitz product over all entries.

Let the selected pivot be $a_{rc}\ne0$, with row and column degrees $u$ and
$v$.  Only the $(u-1)(v-1)$ pairs in its row-tail by column-neighborhood
product can change.  For every such pair, perform
\begin{equation}
 a_{ij}\leftarrow a_{ij}-a_{ic}a_{rc}^{-1}a_{rj}.
 \label{eq:sparse-update}
\end{equation}
Insert a new incidence when a zero becomes nonzero, and delete an incidence
when cancellation gives zero.  Remove the pivot column entries and then the
pivot row.  Column incidence sets make it unnecessary to scan all active
rows merely to discover the ones having $a_{ic}\ne0$.

Complete row and column pivoting is convenient here.  Two compact arrays
record the current row and column orders, with inverse-position arrays.
Move the selected row and column to their last active positions; each
nontrivial transposition flips a determinant sign.  Their last positions
are then deleted in constant time.  Stable original indices continue to
identify sparse rows and columns, so no global relabeling is needed after
a pivot.

\begin{proposition}[Determinant correctness]
\label{prop:sparse-correct}
For a square matrix over a prime field, the incidence-indexed algorithm
returns its exact determinant.  It returns zero upon finding a zero active
row.  It does not mutate the caller's input rows.
\end{proposition}

\begin{proof}
After the recorded permutations, partition the active matrix as
\[
 \begin{pmatrix}B&b\\c^{\mathsf T}&a\end{pmatrix},\qquad a\ne0.
\]
Subtract $b_i/a$ times its last row from each earlier row.  These elementary
operations preserve determinant and yield a block triangular matrix with
diagonal blocks $B-ba^{-1}c^{\mathsf T}$ and $a$.  Its determinant is
$a\det(B-ba^{-1}c^{\mathsf T})$.  Equation~\eqref{eq:sparse-update}
computes exactly this Schur complement.  Entries outside the indicated
Cartesian product do not change.  The row and column transpositions are
accounted for by their signs.  Induction on active dimension proves the
result, with empty determinant one.  A zero active row makes the remaining
determinant zero and therefore the original determinant zero.  Copying and
normalizing the input maps before elimination proves the last assertion.
\end{proof}

The generic function takes primality as a precondition.  It checks the
basic integer and index contracts but does not run a primality test.  The
filter's fixed modulus satisfies the precondition.  Supplying a composite
modulus to this field algorithm is outside its interface contract.

\subsection{An explicit bound in terms of Schur work}

For the pivot sequence actually chosen, define
\begin{equation}
 F=\sum_{\text{pivots }(r,c)}(u_{rc}-1)(v_{rc}-1).
 \label{eq:sparse-fill-work}
\end{equation}
Here $m$ is matrix dimension, $E$ counts all stored input entries inspected
before removing modular zeros, and $Z_{\max}$ is the maximum active nonzero
count.  The parameter $F$ counts attempted update positions, including updates that cancel to
zero and repeated updates to the same position at different pivots.  It is
more informative than counting only newly created nonzeros.

\begin{theorem}[Sparse elimination resource bound]
\label{thm:sparse-work}
Under the dictionary lookup model, elimination uses $O(m+E+F)$ field
operations, together with at most $m$ field inversions.  Its incidence,
ordering, and priority work is
\[
 O\bigl((m+E+F)\log(m+1)\bigr).
\]
Its logical sparse storage, including retained caller-input references, is
$O(m+E+Z_{\max})$ field entries and indices.  A conservative bound allowing
Python dictionaries and sets to retain allocated capacity after deletion is
$O(m+E+F)$.  Every created fill entry is
charged to one of the $F$ update positions.  No favorable bound on $F$ or
$Z_{\max}$ follows solely from the initial $O(n)$ Alexander sparsity.
\end{theorem}

\begin{proof}
Initial normalization and incidence construction inspect $E$ entries.
Every update in \eqref{eq:sparse-update} costs a bounded number of field
operations.  Forming multipliers for rows meeting a pivot column is charged
to those incidences.  Deleting incidences and scanning pivot rows can also
be charged to entries: an entry is either present initially or was created
by a preceding update, and it can be deleted only once before another
creation.  Their total number is at most $E+F$.  Thus even pivots with an
empty row tail, whose product term is zero, have their incidence work
accounted for.

Each affected row inserts a refreshed degree into the priority queue.
Stale records are discarded on extraction.  A queue record is inserted or
discarded only a constant number of times, and the insertion count is
$O(m+E+F)$.  The implementation rebuilds the heap whenever stale records
make it larger than a constant multiple of the active row count, apart
from a fixed small allowance.  A rebuild costs linear time in active rows
and is amortized against records accumulated or rows removed since the
previous rebuild.  Heap operations and stable sorting of pivot-column
incidences therefore fit the stated logarithmic overhead.  The compact
permutation arrays use linear space.

The active matrix and its dual incidence representation use $O(Z_{\max})$
logical entries; the bounded heap uses $O(m)$ records.  The source row list retains
the original input, giving the additional $E$ allowance.  A nonzero first
created during elimination arises in an explicitly visited update position,
so the fill-creation count is at most $F$.  Charging retained allocation to
all entries ever inserted yields the conservative concrete storage bound.
\end{proof}

For a variable modulus, multiply the arithmetic-operation bound by the
appropriate cost of operations on $O(\log p)$-bit integers and include the
cost of modular inversion.  The fixed 61-bit filter avoids coefficient
growth.  This is one reason the modular filter is useful even when an exact
polynomial computation will eventually be needed.

The bound permits cubic work: all active degrees can be proportional to the
active dimension.  Conversely, if a matrix family has $E=O(m)$ and the chosen
pivots satisfy $F=O(m)$, this implementation uses nearly linear work in the
stated model.  This is a conditional consequence of the measured parameter,
not a claimed structural theorem for every knot-derived Alexander matrix.

\subsection{Validation and practical interpretation}

Independent determinant checks use the Leibniz permutation formula on 450
small matrices over three fields, followed by 420 comparisons with the
archived dense elimination over larger dimensions and four fields.  They
include singular matrices, zero rows, characteristic two, complete-pivot
sign changes, and permutation matrices through dimension 1,000.  The latter
require no Schur updates, testing the distinction between pivot-incidence
work and $F$.  Diagram-level checks compare exact witness dictionaries on
120 generated braid knots and selected torus knots, and verify input
preservation and resource callbacks.

The larger torus-knot examples in Section~\ref{sec:experiments} exhibit small
Schur neighborhoods and substantial measured improvements.  Small examples
can be slower because maintaining incidence sets and priorities costs more
than traversing short dense rows.  The archive retains those regressions.
No empirical line fit is offered as a proof that the observed fill pattern
persists on other families.
